%%%%%%%%%%%%%%%%%%%%%%%%%%%%%%%%%%%%%%%%%%%%%%%%%%%%%%%%%%%%%%%%%%%%%%%%%%%%%%%
% Artha — A Domain Ontology-Driven Agentic Framework for LLM-Based
% Personal Finance Reasoning
%
% Target: arXiv preprint (standalone) / IEEE Big Data 2026 (conference)
% Required: IEEEtran.cls (IEEE conference template)
%%%%%%%%%%%%%%%%%%%%%%%%%%%%%%%%%%%%%%%%%%%%%%%%%%%%%%%%%%%%%%%%%%%%%%%%%%%%%%%
\documentclass[conference]{IEEEtran}
\IEEEoverridecommandlockouts

\usepackage{cite}
\usepackage{amsmath,amssymb,amsfonts}
\usepackage{algorithmic}
\usepackage{graphicx}
\usepackage{textcomp}
\usepackage{xcolor}
\usepackage{booktabs}
\usepackage{url}
\usepackage{hyperref}
\hypersetup{colorlinks=true, linkcolor=blue, urlcolor=blue, citecolor=blue}
\def\BibTeX{{\rm B\kern-.05em{\sc i\kern-.025em b}\kern-.08em
    T\kern-.1667em\lower.7ex\hbox{E}\kern-.125emX}}

\begin{document}

\title{Artha: A Domain Ontology-Driven Agentic Framework for
LLM-Based Personal Finance Reasoning}

\author{\IEEEauthorblockN{Tejas Katika}
\IEEEauthorblockA{\textit{Department of Computer Science} \\
\textit{University of North Texas}\\
Denton, TX, USA \\
tejashwar1029@gmail.com}}

\maketitle

%=============================================================================
\begin{abstract}
Large language models (LLMs) are well-suited to conversational personal
finance applications but struggle when reasoning directly over raw
banking transactions, which lack typed semantics and structured
relationships. We present \textbf{Artha}, a personal finance framework
that enriches raw transactions into typed ontology objects
(merchant profiles, spending categories, recurring bills, anomaly
flags, budgets, and goals) before LLM reasoning, and exposes them to
a Claude Sonnet~4.6 agent through 15 typed tools.
We evaluate Artha on a 60-query benchmark spanning 6
financial-reasoning categories and 10 user archetypes, using a
held-out Claude Opus~4.7~\cite{claude_opus} instance as an
LLM-as-judge~\cite{mtbench} over a 4-dimensional rubric. Artha's full configuration
achieves a \textbf{75.0\% pass rate} (3.67/5 mean score) versus a
\textbf{48.3\% raw-baseline} pass rate (Condition D, $\Delta=-26.7$\%),
confirming the value added by ontology-driven enrichment and the typed
tool layer. An ablation study isolates the contribution of each
component: the ontology tool layer drives most of the lift
($\Delta=-25.0$\% when bypassed; Condition B), while enrichment data
alone contributes less because the system gracefully re-computes
missing enrichments on demand ($\Delta=-5.0$\%; Condition C). We report
honest limitations, including single-seed per archetype, post-hoc
detector threshold tuning, same-provider LLM judge, and synthetic
evaluation data. Code, data schema, and evaluation harness are
released for reproducibility.
\end{abstract}

\begin{IEEEkeywords}
LLM agents, domain ontology, personal finance, tool-augmented LLMs,
LLM-as-judge, ablation study
\end{IEEEkeywords}

%=============================================================================
\section{Introduction}

General-purpose LLMs are increasingly used as conversational
interfaces for personal financial management. However, reasoning
directly over raw transactional data (CSV rows or JSON records
containing merchant strings, amounts, dates, and type codes) has
well-known failure modes: inconsistent category attribution,
unstable merchant canonicalization, shallow analytics beyond
simple sums, and hallucinated recurring charges. These failures
matter because personal finance is a high-stakes numerical domain
where confident-but-wrong outputs erode user trust.

We hypothesize that much of this gap is not a language-understanding
problem but a \emph{knowledge-representation} problem. Raw
transactions lack typed semantics, cross-transaction structure, and
temporal aggregates that are essential for questions such as
\emph{``what recurring subscriptions do I have?''} or
\emph{``am I saving enough for my goals?''}. A thin typed ontology
layer, populated by deterministic enrichment and surfaced through
typed agent tools, can recover most of this structure while remaining
interpretable and auditable.

This paper presents \textbf{Artha}, an end-to-end personal finance
agent built on this hypothesis, and evaluates it under a controlled
ablation study. Our contributions are:

\begin{itemize}
  \item \textbf{A domain ontology} for personal finance
    comprising 9 core object types and 7 primary relationships
    (Section~\ref{sec:ontology}), with a dual-source enrichment
    pipeline (rule-based merchant classification plus upstream-label
    fallback) that reaches 97\% typed categorization coverage on
    synthetic transactions.
  \item \textbf{A 15-tool typed agent layer} that queries the ontology
    rather than raw tables (Section~\ref{sec:tools}), allowing an LLM
    to produce numerically grounded responses to natural-language
    finance questions.
  \item \textbf{A reproducible 60-query benchmark} and LLM-as-judge
    harness covering 6 reasoning categories across 10 user archetypes,
    with documented rubric and release of raw responses and scores
    (Section~\ref{sec:eval}).
  \item \textbf{Ablation results} isolating the contributions of
    enrichment data versus the typed tool layer
    (Section~\ref{sec:ablation}), showing that the tool layer drives
    most of the observed lift.
\end{itemize}

Throughout, we report honest methodological
compromises---reference-date anchoring, post-hoc detector
threshold tuning on synthetic data, same-provider judge, and
single-user-per-archetype evaluation---and defer the stronger
validation steps (human Cohen's $\kappa$, cross-family judge,
multi-user per archetype) to future work.

%=============================================================================
\section{Related Work}
\label{sec:related}

\paragraph{Tool-augmented LLMs and agents.}
Tool use~\cite{toolformer,react} and structured tool-calling
APIs~\cite{anthropic_tools} have become standard for enabling LLMs
to retrieve and manipulate external state. Artha follows this
pattern: a Claude Sonnet~4.6~\cite{claude_sonnet} model is given a
strict tool schema and
orchestrates calls to typed finance tools via a bounded
reasoning loop (max 8 iterations per query).

\paragraph{Ontology-guided reasoning.}
A long line of work has used ontologies to ground language-model
reasoning in typed, interlinked objects~\cite{ontology_survey}.
Most such work has focused on knowledge-graph question answering;
ours targets transactional reasoning, where the ``ontology'' is
narrow, financial, and populated by rule- and ML-based enrichment
rather than curated triples.

\paragraph{LLM-as-judge.}
MT-Bench~\cite{mtbench} and subsequent works validated the use of
strong LLMs (e.g., GPT-4, Claude Opus) as scalable judges of
open-ended LLM outputs. We adopt this methodology, using a held-out
Claude Opus~4.7 judge with an explicit 4-dimensional rubric, and
discuss its same-provider limitation in Section~\ref{sec:limitations}.

\paragraph{LLM-based personal finance assistants.}
Commercial products (Rocket~Money, Cleo, Copilot Money) have
popularized LLM-backed finance assistants, but published evaluations
are scarce. Our 60-query benchmark and open ablation harness aim to
provide a reproducible reference point.

%=============================================================================
\section{The Artha Framework}
\label{sec:framework}

Artha is implemented as a Java 21 / Spring Boot~3.2 service backed by
PostgreSQL~15, with a Claude Sonnet~4.6 agent orchestrating calls to
typed ontology tools. Figure~\ref{fig:arch} summarizes the four
layers, from raw bank transactions to the user-facing agent response.

\begin{figure}[t]
\centering
\fbox{\parbox{0.92\columnwidth}{\centering
\textbf{Layer 1} -- Raw transactions, bank accounts \\
$\downarrow$ \\
\textbf{Layer 2} -- Ontology enrichment \\
\textit{(merchant classification, category resolution, anomaly detection,
recurring-bill detection)} \\
$\downarrow$ \\
\textbf{Layer 3} -- Typed ontology objects \\
\textit{(9 object types, 7 primary relationships)} \\
$\downarrow$ \\
\textbf{Layer 4} -- 15 agent tools over typed objects \\
$\downarrow$ \\
\textbf{Claude Sonnet 4.6} agent loop \\
(system prompt, tool-call iterations, response synthesis)
}}
\caption{Architecture of the Artha framework. Enrichment is
deterministic and idempotent; the LLM only reads the enriched objects
through typed tool calls.}
\label{fig:arch}
\end{figure}

%-----------------------------------------------------------------------------
\subsection{Domain Ontology}
\label{sec:ontology}

The ontology contains nine core object types. Five derive from raw
banking records, and four are produced or refined by enrichment:

\begin{itemize}
  \item \textbf{Transaction}: raw debit or credit record with amount,
    post-date, merchant-name string, and a JSON metadata field
    carrying the upstream canonical category (as provided, for example,
    by the Plaid Transactions API~\cite{plaid}).
  \item \textbf{TransactionEnrichment}: 1:1 projection of a
    transaction carrying derived fields
    (\texttt{spending\_category\_id}, \texttt{merchant\_profile\_id},
    \texttt{is\_anomaly}, \texttt{anomaly\_reason},
    \texttt{canonical\_merchant\_name},
    \texttt{budget\_utilization\_pct}, \texttt{enrichment\_source}).
  \item \textbf{MerchantProfile}: canonical merchant identity
    (\emph{Netflix}, \emph{T-Mobile}, \ldots) with
    \texttt{is\_recurring} flag and links to a \textbf{MerchantType}.
  \item \textbf{MerchantType}: 19 seed categories
    (\emph{Restaurant}, \emph{Utilities}, \emph{Streaming}, \ldots).
  \item \textbf{SpendingCategory}: per-user category instance
    linked to a merchant type (e.g., a user's ``Restaurant''
    spending bucket).
  \item \textbf{ClassificationRule}: pattern-based rule mapping
    normalized merchant strings to merchant profiles
    (\texttt{pattern}, \texttt{pattern\_type}, \texttt{priority}).
  \item \textbf{RecurringBill}: detected recurring charge with
    \texttt{expected\_amount}, \texttt{billing\_cycle},
    \texttt{next\_expected\_date}, and \texttt{confidence\_score}.
  \item \textbf{Budget}: per-user, per-category monthly limit with
    \texttt{effective\_from} / \texttt{effective\_to} validity
    window.
  \item \textbf{FinancialGoal}: target amount, current amount,
    monthly contribution, target date, and status
    (\emph{ACTIVE}, \emph{ACHIEVED}, \emph{PAUSED}).
\end{itemize}

The primary relationships between these objects (Table~\ref{tab:rel})
give the LLM a traversable, typed graph rather than a flat table.

\begin{table}[t]
\centering
\caption{Primary ontology relationships.}
\label{tab:rel}
\begin{tabular}{@{}ll@{}}
\toprule
Source $\rightarrow$ Target & Cardinality \\
\midrule
Transaction $\rightarrow$ TransactionEnrichment & 1:1 \\
TransactionEnrichment $\rightarrow$ SpendingCategory & N:1 \\
TransactionEnrichment $\rightarrow$ MerchantProfile & N:1 \\
MerchantProfile $\rightarrow$ MerchantType & N:1 \\
ClassificationRule $\rightarrow$ MerchantProfile & N:1 \\
Budget $\rightarrow$ SpendingCategory & N:1 \\
RecurringBill $\rightarrow$ MerchantProfile & N:1 \\
\bottomrule
\end{tabular}
\end{table}

%-----------------------------------------------------------------------------
\subsection{Dual-Source Enrichment Pipeline}
\label{sec:enrichment}

Each incoming transaction passes through a four-step enrichment
pipeline, producing a \texttt{TransactionEnrichment}:

\begin{enumerate}
  \item \textbf{Merchant classification}: normalize the merchant
    string (uppercase, trim) and match against the priority-ordered
    \texttt{ClassificationRule} set (\emph{EXACT}, \emph{PREFIX},
    \emph{CONTAINS}, \emph{REGEX} patterns). A successful match
    attaches a \texttt{MerchantProfile} and a confidence weight
    $\in [0.70, 0.99]$.
  \item \textbf{Category resolution}: three-tier fallback.
    (1)~If classification matched, use the merchant profile's
    \texttt{MerchantType} to \texttt{findOrCreate} the user's
    \texttt{SpendingCategory}. (2)~Otherwise, read the upstream
    canonical category from \texttt{transaction.metadata['category']}
    (simulating a bank-API--provided label in production) and resolve
    it to the user's category by name. (3)~Otherwise, fall back to
    \emph{Income} for credits / \emph{Other} for debits.
  \item \textbf{Budget utilization}: if the transaction falls inside
    an active user budget for the resolved category, compute and
    record month-to-date utilization percent.
  \item \textbf{Anomaly detection}: a per-merchant statistical
    detector flags a transaction as anomalous when its amount is
    more than $Z = 2.5$ standard deviations above the merchant's
    30-day rolling mean (minimum sample size = 3 prior
    transactions). A fallback rule flags amounts $> 3\times$ the
    merchant profile's \texttt{typical\_amount\_max}.
\end{enumerate}

The enrichment is recorded with an
\texttt{enrichment\_source} label (\emph{RULES}, \emph{METADATA},
\emph{FALLBACK}). In our 10-archetype evaluation set, 40.4\% of
enrichments are produced by rule matches against 23 seed merchant
rules and 59.6\% by the upstream-metadata fallback; the type-based
default path is not triggered on this dataset.

\textbf{Recurring-bill detection} is a separate sweep over each
user's last 180~days of debits. A merchant with at least four
occurrences and an average inter-transaction interval in $[5, 40]$
days is emitted as a \emph{monthly} \texttt{RecurringBill}, and
$[355, 380]$ days as \emph{annual}. We disabled a stricter
amount-consistency filter because the synthetic generator used for
evaluation produces high-variance amounts per merchant; this choice
is described in Section~\ref{sec:limitations}.

%-----------------------------------------------------------------------------
\subsection{Agent Tool Layer}
\label{sec:tools}

The LLM never executes SQL. It is given a strict JSON-schema view of
15 typed tools, grouped into four categories:

\begin{itemize}
  \item \textbf{Data}: \texttt{get\_spending\_summary},
    \texttt{get\_transactions}, \texttt{get\_income\_vs\_spend}.
  \item \textbf{Ontology}: \texttt{get\_anomalies},
    \texttt{get\_subscriptions}, \texttt{get\_goal\_progress},
    \texttt{get\_financial\_health},
    \texttt{get\_category\_insights}.
  \item \textbf{Planning}: \texttt{get\_what\_if\_analysis},
    \texttt{get\_investment\_projection},
    \texttt{get\_goal\_projection}, \texttt{get\_goal\_impact}.
  \item \textbf{Analytics}: \texttt{get\_budget\_status},
    \texttt{get\_behavioral\_patterns},
    \texttt{get\_financial\_education}.
\end{itemize}

Each tool returns structured JSON (dollar totals, category breakdowns,
recurring-bill records, ranked insights). The agent is configured
with \texttt{max\_tokens=4096}, \texttt{temperature=0.3}, and a
hard cap of 8 tool-call iterations per user query.

Crucially, tool results are carried into the LLM's context as
text-encoded JSON, not as free-form prose: the LLM performs
summarization and reasoning but does not invent numerical facts.

%=============================================================================
\section{Evaluation Methodology}
\label{sec:eval}

%-----------------------------------------------------------------------------
\subsection{Synthetic Dataset}

We evaluate on a fixed synthetic dataset of 509 users and
128{,}096 debit/credit transactions spanning 2024-01-01 to
2026-02-28. Users are distributed across 10 archetypes:
\emph{high\_earner}, \emph{lifestyle\_overspender},
\emph{high\_debt}, \emph{paycheck\_to\_paycheck}, \emph{freelancer},
\emph{retired\_fixed}, \emph{dual\_income}, \emph{recent\_grad},
\emph{small\_biz\_owner}, \emph{gig\_worker}. Each archetype has 50
users; for this evaluation we use the first (seed) user per
archetype---a limitation discussed in Section~\ref{sec:limitations}.

To keep the synthetic data temporally coherent with the Artha
tools' ``now'' semantics, we anchor the evaluation \emph{reference
date} at \textbf{2024-12-31} (the end of the primary data window
for the 500 archetype users). All tool queries that reference
``this month'' or ``last 30 days'' are evaluated against this
reference date instead of wall-clock time. This substitution is
implemented through a single \texttt{ReferenceDateProvider} bean
injected into all time-dependent tools, and is disabled (falls
back to wall-clock) in production.

%-----------------------------------------------------------------------------
\subsection{Benchmark Queries}

The benchmark consists of 60 natural-language queries structured as
6 \emph{financial-reasoning categories} $\times$ 10 \emph{user
archetypes}. The categories are:

\begin{itemize}
  \item \textbf{spending\_summary}:
    \emph{``How much did I spend this month? Break it down by
    category.''}
  \item \textbf{financial\_health}:
    \emph{``How are my finances looking overall? Am I saving
    enough?''}
  \item \textbf{goal\_tracking}:
    \emph{``Am I on track to meet my financial goals?''}
  \item \textbf{anomaly\_detection}:
    \emph{``Have there been any unusual or suspicious charges
    recently?''}
  \item \textbf{subscription\_audit}:
    \emph{``What recurring charges and subscriptions do I have?''}
  \item \textbf{behavioral\_analysis}:
    \emph{``What are my worst spending habits? Where am I wasting
    money?''}
\end{itemize}

%-----------------------------------------------------------------------------
\subsection{LLM-as-Judge Rubric}
\label{sec:rubric}

Each agent response is scored by a held-out Claude Opus~4.7
instance on four dimensions on a 1--5 scale (judge temperature
omitted; Opus~4.7 ignores the parameter):

\begin{itemize}
  \item \textbf{DA}~(Data Accuracy): Does the response cite specific,
    correct dollar amounts, merchant names, dates, or percentages
    from the user's actual financial data?
  \item \textbf{IQ}~(Insight Quality): Is the financial reasoning
    sound and archetype-specific?
  \item \textbf{AC}~(Actionability): Does the response provide
    concrete, prioritized next steps?
  \item \textbf{CO}~(Completeness): Does the response surface all
    important financial signals present in the context?
\end{itemize}

A response \emph{passes} if its mean dimension score is
$\geq 3.5/5$. Both the agent (Claude Sonnet~4.6) and judge
(Claude Opus~4.7) are from the same model provider---a limitation
discussed in Section~\ref{sec:limitations}.

%=============================================================================
\section{Results}
\label{sec:results}

%-----------------------------------------------------------------------------
\subsection{Condition A: Full System}

Under the full configuration, Artha passes \textbf{45 of 60} queries
(\textbf{75.0\%}) with a mean score of \textbf{3.67/5}. Mean query
latency is 12.3 seconds (agent + tool calls; judging not
included). Table~\ref{tab:a_cat} gives the per-category breakdown.

\begin{table}[t]
\centering
\caption{Condition A (full system): per-category results.}
\label{tab:a_cat}
\begin{tabular}{@{}lrr@{}}
\toprule
Category & Pass \% & Mean \\
\midrule
spending\_summary     & 100.0\% & 3.85 \\
behavioral\_analysis  &  90.0\% & 4.30 \\
subscription\_audit   &  90.0\% & 3.80 \\
financial\_health     &  80.0\% & 3.80 \\
goal\_tracking        &  60.0\% & 3.38 \\
anomaly\_detection    &  30.0\% & 2.88 \\
\midrule
\textbf{Overall (n=60)} & \textbf{75.0}\% & \textbf{3.67} \\
\bottomrule
\end{tabular}
\end{table}

Two categories---\texttt{anomaly\_detection} and
\texttt{goal\_tracking}---account for most failures. Manual
inspection confirms that these are largely \emph{truthful empty
results}: the archetype seed users have on average 0.5 statistically
flagged anomalies in the evaluation window, so the tool correctly
returns ``no anomalies found.'' The judge penalizes these responses
for lack of concrete next steps even when there is nothing to flag,
reflecting a well-known difficulty of rubric-based scoring on
correct-but-unhelpful answers.

Table~\ref{tab:a_arch} gives the per-archetype view. Pass rates
range from 50.0\% (\emph{high\_debt}) to 100.0\% (\emph{gig\_worker}),
with archetype-specific variation driven primarily by whether the
archetype's seed user has any anomalies or is data-sparse in
behavioral-pattern detection.

\begin{table}[t]
\centering
\caption{Condition A: per-archetype results.}
\label{tab:a_arch}
\begin{tabular}{@{}lrr@{}}
\toprule
Archetype & Pass \% & Mean \\
\midrule
gig\_worker              & 100.0\% & 3.71 \\
paycheck\_to\_paycheck   &  83.3\% & 3.71 \\
freelancer               &  83.3\% & 3.54 \\
retired\_fixed           &  83.3\% & 3.42 \\
dual\_income             &  83.3\% & 3.71 \\
high\_earner             &  66.7\% & 3.88 \\
lifestyle\_overspender   &  66.7\% & 3.88 \\
recent\_grad             &  66.7\% & 3.71 \\
small\_biz\_owner        &  66.7\% & 3.67 \\
high\_debt               &  50.0\% & 3.46 \\
\bottomrule
\end{tabular}
\end{table}

%-----------------------------------------------------------------------------
\subsection{Ablation Study}
\label{sec:ablation}

To isolate the contribution of each framework component, we run
three ablation conditions alongside the full configuration:

\begin{itemize}
  \item \textbf{A -- Full}: enrichment data populated, ontology
    tools enabled.
  \item \textbf{B -- Tools OFF}: enrichment data populated, but
    the four ontology-aware tools
    (\texttt{get\_anomalies}, \texttt{get\_subscriptions},
    \texttt{get\_category\_insights},
    \texttt{get\_financial\_health})
    are switched to a degraded mode that bypasses enrichment joins
    and returns baseline responses.
  \item \textbf{C -- Enrichment OFF}: all category, profile, and
    anomaly fields on \texttt{TransactionEnrichment} are nulled,
    and \texttt{RecurringBill} rows are removed. Tools remain
    fully enabled.
  \item \textbf{D -- Both OFF}: C plus B (raw baseline).
\end{itemize}

Results are in Table~\ref{tab:ablation}. For Conditions~B and~D we
applied a 2-second inter-query pause to stay below Anthropic tool-use
rate limits; all runs are otherwise identical and use the same 60
benchmark queries.

\begin{table}[t]
\centering
\caption{Ablation study (Table~5 of the paper).
$\Delta$ is the change in overall pass rate relative to Condition A.}
\label{tab:ablation}
\begin{tabular}{@{}llrrr@{}}
\toprule
Cond. & Enrichment / Tools & Pass~\% & Mean & $\Delta$ \\
\midrule
A & $\checkmark$ / $\checkmark$ & 75.0\% & 3.67 & --- \\
B & $\checkmark$ / $\times$     & 50.0\% & 2.97 & $-25.0$\% \\
C & $\times$ / $\checkmark$     & 70.0\% & 3.56 & $-5.0$\%  \\
D & $\times$ / $\times$         & 48.3\% & 2.86 & $-26.7$\% \\
\bottomrule
\end{tabular}
\end{table}

Three findings stand out.

\paragraph{The full system improves over the raw baseline by
26.7~percentage points.} A~$\rightarrow$~D isolates the combined
contribution of ontology-driven enrichment and typed tools.

\paragraph{The ontology tool layer drives most of the lift
($-25.0$\%).} A~$\rightarrow$~B removes only the tool layer while
leaving enrichment data in place. Pass rate collapses almost as
far as the raw baseline, demonstrating that \emph{typed-object query
design matters more than the presence of enrichment data alone}.

\paragraph{Enrichment data alone contributes only $-5.0$\% due to
graceful degradation.} A~$\rightarrow$~C leaves the tool paths in
place but removes their inputs. The modest drop is explained by
an auto-recovery behavior we deliberately preserved:
\texttt{get\_subscriptions} re-runs the subscription detector when
\texttt{recurring\_bills} is empty, and other ontology tools retain
enough fallback logic to produce useful responses. This is
\emph{intentional system robustness}, not a flaw in the
ablation---and we surface it as a legitimate property of the
framework rather than a confound.

Table~\ref{tab:per_cat_ablation} gives the per-category pass rates
across all four conditions. \texttt{anomaly\_detection} and
\texttt{subscription\_audit} are the two categories that collapse
to 0\% under B (tools off), showing that these reasoning modes are
fundamentally typed-tool-mediated. \texttt{spending\_summary} and
\texttt{behavioral\_analysis} remain robust across conditions
because they read from \texttt{transactions.metadata} directly and
do not depend on the enrichment ontology path.

\begin{table}[t]
\centering
\caption{Per-category pass-rate comparison across conditions.}
\label{tab:per_cat_ablation}
\begin{tabular}{@{}lrrrr@{}}
\toprule
Category & A & B & C & D \\
\midrule
spending\_summary    & 100\% & 70\%  & 80\%  & 60\% \\
financial\_health    & 80\%  & 70\%  & 70\%  & 80\% \\
goal\_tracking       & 60\%  & 60\%  & 70\%  & 60\% \\
anomaly\_detection   & 30\%  & 0\%   & 0\%   & 0\%  \\
subscription\_audit  & 90\%  & 0\%   & 100\% & 0\%  \\
behavioral\_analysis & 90\%  & 100\% & 100\% & 90\% \\
\bottomrule
\end{tabular}
\end{table}

%-----------------------------------------------------------------------------
\subsection{Qualitative Example}

The following excerpt is the agent's response under Condition A
to \emph{``What recurring charges and subscriptions do I have?''}
for the \emph{high\_earner} archetype (user seed
\texttt{aa000000\ldots}, judge score 4.0/5):

\begin{quote}\small\itshape
``Here is a breakdown of your 10 detected recurring charges, totalling
\$3{,}772.42/month (or \$45{,}269/year):

1. OpenTable \$432.39 (Jan 23); 2. Tesla \$427.12 (Jan 6);
3. Fidelity \$422.78 (Dec 31); 4. Whole Foods \$412.71 (Jan 11);
5. Equinox \$385.94 (Dec 27); 6. Delta Airlines \$362.82 (Jan 15);
7. Vanguard \$361.89 (Jan 4); 8. Sweetgreen \$355.64 (Jan 24);
\ldots''
\end{quote}

Every dollar figure and date in this response is produced by the
\texttt{SubscriptionDetector} and surfaced through
\texttt{get\_subscriptions}; the LLM's role is to translate the
structured result into readable natural language without modifying
the numbers. Under Condition~B, the same query returns:

\begin{quote}\small\itshape
``Unfortunately, I was not able to retrieve your subscription data
at the moment. The subscription detection service is currently
unavailable (the recurring-bill detection pipeline is bypassed).''
\end{quote}

The judge score drops from 4.0 to 1.0, illustrating why
\texttt{subscription\_audit} collapses from 90\% to 0\% between
A and B.

%=============================================================================
\section{Discussion}

\subsection{Multi-Source Enrichment Is a Feature, Not a Fallback}

A key implementation decision in Artha is to treat upstream-provided
categorization (e.g., the category field returned by Plaid or a
similar bank-data aggregator) as a \emph{first-class enrichment
source} alongside rule-based classification. In production, this
matches how real financial data arrives: Plaid-style APIs already
provide a canonical category for each transaction, and it would be
wasteful for a thin rule engine to re-derive that label when a
high-quality label is already present.

We report the per-source coverage transparently
(Section~\ref{sec:enrichment}): 40.4\% of the enriched transactions
in our evaluation set were labeled by the rule engine and 59.6\% by
upstream metadata. A reader interested in a ``pure rule-engine''
number should read the 40.4\% figure, not the near-100\% combined
coverage.

\subsection{Graceful Degradation as an Architectural Property}

The modest A~$\rightarrow$~C gap is a finding, not a failure. The
system is designed to be robust to missing enrichments: when
\texttt{recurring\_bills} is empty, the subscription tool re-runs
detection; when \texttt{spending\_category\_id} is \texttt{NULL},
tools that can fall back to \texttt{metadata[category]} do so.
This keeps the user experience consistent when the enrichment
pipeline is cold, under migration, or partially unavailable---a
property that matters for a 24/7 production deployment.

%=============================================================================
\section{Limitations}
\label{sec:limitations}

We explicitly document the following compromises to allow the reader
to calibrate the results:

\begin{enumerate}
  \item \textbf{Single seed user per archetype.} Each archetype is
    evaluated with one user (the first of 50). This bounds
    generalization across archetype variance; we plan to expand to
    $n \geq 5$ users per archetype in future work.
  \item \textbf{Synthetic evaluation data.} All transactions are
    generated by a controlled generator. We rely on the upstream
    metadata category field, which in real systems would be bank-provided.
  \item \textbf{Post-hoc detector threshold tuning.} The recurring-bill
    detector's monthly-interval window was widened from the
    original $[25, 35]$ days to $[5, 40]$ days and the amount-variance
    filter was disabled to accommodate the generator's intentionally
    noisy amounts. A production deployment against real bank data
    would re-enable the amount-variance filter with an appropriate
    threshold.
  \item \textbf{Reference-date substitution.} Tools resolve ``now''
    to 2024-12-31 so that the dataset's finite window intersects the
    tools' rolling query windows. The override is a single injected
    bean and is disabled in production.
  \item \textbf{Same-provider judge.} Both agent (Claude Sonnet~4.6)
    and judge (Claude Opus~4.7) are from Anthropic. A cross-family
    second judge is valuable for bias mitigation and is deferred to
    future work.
  \item \textbf{No human-annotator validation.} We did not compute
    Cohen's $\kappa$~\cite{cohens_kappa} between human annotators and
    the LLM judge for this version. The scoring rubric
    (\texttt{SCORING\_RUBRIC.md}) is released to enable external
    validation.
  \item \textbf{Seeded goals and budgets.} To establish a consistent
    starting state across archetypes, we seeded two
    \texttt{FinancialGoal}s and three \texttt{Budget}s per user
    uniformly. These are evaluation-setup data, not outputs of the
    Artha system; the paper's claims concern the enrichment and tool
    layers, not goal/budget personalization.
  \item \textbf{Limited rule set.} Our rule corpus contains 23
    classification rules (8 seeded plus 15 added for the top
    merchants by transaction volume). This is intentionally small
    for a research prototype; production systems would scale to
    hundreds or thousands of rules.
\end{enumerate}

%=============================================================================
\section{Conclusion}

We presented Artha, a domain-ontology-driven agentic framework
for personal finance reasoning, and evaluated it on a 60-query
benchmark with full ablation coverage. Artha's full configuration
reaches 75.0\% pass rate (3.67/5 mean score), and the ablation
isolates the contribution of the typed-tool layer as the dominant
source of lift ($-25$\% when disabled), ahead of the enrichment
data itself ($-5$\%). The framework degrades gracefully when
enrichment is missing thanks to the tools' ability to re-compute
on demand.

Beyond the quantitative results, Artha's value proposition is
\emph{structural}: constraining the LLM to reason over typed
ontology objects through a fixed-surface tool API yields
numerically grounded responses that are auditable, reproducible,
and straightforward to extend.

Future work includes expanding evaluation to multiple users per
archetype, adding a cross-family LLM judge, computing Cohen's
$\kappa$ against human annotators, and replicating the evaluation
on a small real-bank dataset under appropriate privacy controls.

%=============================================================================
% Inline bibliography so the arXiv submission is self-contained
% (arXiv re-runs LaTeX but not BibTeX, so a .bbl would be required
% otherwise; inlining avoids that dependency entirely).
\begin{thebibliography}{10}

\bibitem{toolformer}
T.~Schick, J.~Dwivedi-Yu, R.~Dess{\`\i}, R.~Raileanu, M.~Lomeli,
L.~Zettlemoyer, N.~Cancedda, and T.~Scialom,
``{Toolformer}: Language models can teach themselves to use tools,''
in \emph{Advances in Neural Information Processing Systems (NeurIPS)},
2023.

\bibitem{react}
S.~Yao, J.~Zhao, D.~Yu, N.~Du, I.~Shafran, K.~Narasimhan, and Y.~Cao,
``{ReAct}: Synergizing reasoning and acting in language models,''
in \emph{International Conference on Learning Representations (ICLR)},
2023.

\bibitem{anthropic_tools}
Anthropic, ``Tool use with {Claude},'' Anthropic API documentation,
2024. [Online]. Available:
\url{https://docs.anthropic.com/en/docs/build-with-claude/tool-use}

\bibitem{ontology_survey}
D.~Diefenbach, V.~Lopez, K.~Singh, and P.~Maret,
``A survey on ontology-based knowledge graph question answering,''
\emph{Knowledge and Information Systems}, vol.~55, pp. 529--569, 2018.

\bibitem{mtbench}
L.~Zheng, W.-L. Chiang, Y.~Sheng, S.~Zhuang, Z.~Wu, Y.~Zhuang, Z.~Lin,
Z.~Li, D.~Li, E.~P. Xing, H.~Zhang, J.~E. Gonzalez, and I.~Stoica,
``Judging {LLM}-as-a-judge with {MT-Bench} and chatbot arena,''
in \emph{Advances in Neural Information Processing Systems (NeurIPS)
Datasets and Benchmarks Track}, 2023.

\bibitem{claude_sonnet}
Anthropic, ``{Claude Sonnet 4.6},'' 2026. [Online]. Available:
\url{https://www.anthropic.com/claude}

\bibitem{claude_opus}
Anthropic, ``{Claude Opus 4.7},'' 2026. [Online]. Available:
\url{https://www.anthropic.com/claude}

\bibitem{plaid}
Plaid, ``Plaid {Transactions API}: category classification,'' 2024.
[Online]. Available:
\url{https://plaid.com/docs/api/products/transactions}

\bibitem{cohens_kappa}
J.~Cohen, ``A coefficient of agreement for nominal scales,''
\emph{Educational and Psychological Measurement}, vol.~20, no.~1,
pp. 37--46, 1960.

\end{thebibliography}

\end{document}
